%\clearpage
%\chapter*{Abstract}
%\addcontentsline{toc}{chapter}{\protect Abstract}

Aerodynamic heating poses a major challenge in hypersonic flight, making thermal design a central consideration in vehicle development. This requires accurate prediction of both the heating imposed by the flow and the resulting temperatures within the vehicle structure. Coupled simulations of the flow and structural heat conduction throughout the flight trajectory are computationally demanding, limiting the speed of design iterations. Wall functions offer a way to reduce this cost by relaxing the near-wall mesh requirements of Reynolds-averaged Navier--Stokes simulations. This work investigates existing analytical wall functions, implements them within a computational framework, and systematically assesses their accuracy in predicting aerothermal heat-flux.
By isolating individual sources of error, the study identifies promising wall-function closures but reveals that their performance is substantially constrained by the underlying turbulence models and conventional meshing practices. Both retain assumptions that are shown to not reliably extend to hypersonic flows over strongly cooled walls. A novel mesh-resolution criterion addresses the latter limitation and challenges the longstanding assumption that wall functions inherently struggle in the buffer layer. The turbulence-model assessment further identifies the assumption of a constant turbulent Prandtl number as a likely source of discrepancies in predicted heat-flux and thermal boundary layer profiles. These findings demonstrate that reliable application of analytical wall functions requires joint assessment of closure accuracy, mesh resolution, and turbulence modelling, and identify priorities for improving the accuracy and computational efficiency of hypersonic aerothermal predictions.

Furthermore, proceeding through this work, various important corrections were discovered and successfully submitted for correction, most notably a correction to the momentum thickness definition in Schlichting's Boundary Layer Theory \cite{SchlichtingGersten2017BoundaryLayer}, a typo in the reported viscosity law constant, hindering reproducibility, in Malik and Anderson's laminar boundary layer paper \cite{MalikAnderson1991RealGas} and an error in the implementation of the compressible Spalart--Allmaras turbulence model in OpenFOAM. For completeness, the accepted proofs and arguments supporting these and other corrections are in the appendices~\ref{app:malik-anderson-sutherland} and~\ref{app:spalart-allmaras-density}.